This final project of the Specialist Course "Lattice Field Theory" (by Prof. Dr. Fernando Romero-López and Victor Granados) aims to do a comparison of the Metropolis and the HMC algorithms applied to scalar $\phi^4$ theory using Markov Chain Monte Carlo methods. The main objective is to compare the performance, efficiency, and physical behavior of the Metropolis and Hybrid Monte Carlo (HMC) algorithms.\\\\
The project combines algoritmic and computational aspects. The Metropolis and HMC algorithms were implemented in Python and Rust. For HMC, two integrators (leapfrog and omelyan2) were used and their effect on energy conservation was studied. \\\\
The simulations were carried out for the two-dimensional $\phi^4$ theory in both the symmetric and broken phases. A range of lattice sizes were considered in order to study errors, autocorrelations and observables as functions of $L$. The topological charge was investigated using periodic and anti-periodic boundary conditions in order to probe non-trivial topological effects such as kink–antikink configurations.\\\\
The main conclusion of this final project is that HMC with omelyan2 integrator in Rust provides the most favorable combination of accuracy, efficiency, and low autocorrelation.